\section{Taylorův polynom}
\begin{motivace}
	Máme-li funkci $f$, která má v bodě $a$ vlastní derivaci, pak je tečna
	$$
	    t(x) = f(a) + f^\prime(a)(x-a).
	$$
	Podíváme-li se na limitu
	$$
 \lim_{x \to a} \frac{f(x)-t(x)}{x-a} = 0 = \lim_{x \to a} \frac{f(x)-f(a)-f^\prime(a)(x-a)}{x-a},
	$$
	zjistíme, že aproximuje chování funkce.
\end{motivace}

\subsection{Základní vlastnosti}
\begin{definice}
Nechť $f$ je funkce, $a \in \mathbb{R}$ a $f^{(n)}(a) \in \mathbb{R}$. Pak polynom
\[
T^{f,a}_n(x) = f(a) + f'(a)(x - a) + \dots + \frac{1}{n!} f^{(n)}(a)(x - a)^n
\]
nazýváme \textbf{Taylorovým polynomem} řádu $n$ funkce $f$ v bodě $a$.
\end{definice}

\begin{poznamka}\,
\begin{enumerate}[(1)]
	\item $\operatorname{st} \left( T_n^{f,a} \right ) \leq n $
	\item $T:=T_n^{f,a}$, pak $T(a) = f(a)$,
	$$T^\prime(a) = f^\prime(a)+ f^{\prime\prime}(a)(x-a)+\dots+ \frac{1}{(n-1)!}f^{(n)}(a)(x-a)^{n-1}=T_{n-1}^{f^\prime,a}(a)=f^\prime(a)$$
	$$
	    T^{\prime\prime}(a) = f^{\prime\prime}(a)\quad \dots\quad T^{(n)}(a)=f^{(n)}(a)
	$$
\end{enumerate}
\end{poznamka}
